\begin{answer}
    
i) To derive the range of $\alpha$ and determine $\Theta^{[t]}$, we need to first define the gradient of our objective function:
\[
\nabla J(\Theta) = \beta\Theta
\]

Therefore the GD Update function is as follows:
\[
\Theta^{[t]} = \Theta^{[t-1]} - \alpha \cdot (\beta\Theta^{[t-1]})
\]
    

    We can also define \(\Theta^{[t+1]}\) in terms of the series of updates starting from \(\Theta^{[t]}\).
\begin{align*}
\Theta^{[t+1]} &= (1-\alpha\beta) \Theta^{[t]} \\
\Theta^{[t]} &= (1-\alpha\beta) \Theta^{[t-1]} \\
\Theta^{[1]} &= (1-\alpha\beta)\Theta^{[0]} \\
\Theta^{[2]} &= (1-\alpha\beta)\Theta^{[1]} = (1-\alpha\beta)^2 \Theta^{[0]} \\
\Theta^{[3]} &= (1-\alpha\beta)\Theta^{[2]} = (1-\alpha\beta)^3 \Theta^{[0]}
\end{align*}

When continuing this series, we can generally say that:
\[
\Theta^{[t]} = (1-\alpha\beta)^t \Theta^{[0]}
\]

We can now calculate the value of \(\Theta^\dagger\) when Gradient Descent converges.
\[
\lim_{t \to \infty} |\Theta^{[t]} - \Theta^\dagger| = 0
\]
\[
\lim_{t \to \infty} |(1-\alpha\beta)^t \cdot \Theta^{[0]} - \Theta^\dagger| = 0
\]
As \(t \to \infty\), the term \((1-\alpha\beta)^t \cdot \Theta^{[0]}\) converges towards zero.
\[
\lim_{t \to \infty} \Theta^{[t]} = 0
\]
\(\Rightarrow\) So it follows \(\Theta^\dagger = 0\).

\vspace{1em}
\textbf{Range of \(\alpha\) so that GD converges.}

For the above-mentioned limit to converge, the following needs to be true:
\[
|1-\alpha\beta| < 1
\]
This is equivalent to:
\[
-1 < 1-\alpha\beta < 1
\]
So:
\begin{itemize}
    \item[i)] \(-1 < 1-\alpha\beta \iff \alpha\beta < 2 \iff \alpha < \frac{2}{\beta}\)
    \item[ii)] \(1-\alpha\beta < 1 \iff -\alpha\beta < 0 \iff \alpha\beta > 0\)
\end{itemize}
It follows that for the range of $\alpha$, we have \(0 < \alpha < \frac{2}{\beta}\).

\vspace{1em}
\textbf{Global Minimum:}
\begin{alignat*}{2}
    \nabla J(\Theta^*) &= 0 \\
    \beta\Theta^*&=0 \hspace{1em} &&\left| : \beta; \beta > 0  \right. \\
    \Theta^* &= 0
  \end{alignat*}

The global minimum \(\Theta^*\) equals \(\Theta^\dagger\).

\vspace{1em}
ii)


\textbf{Minimum Number of Iterations}

For the minimum number of iterations \(T\), the following is true:
\begin{alignat*}{2}
    |\Theta^{[T]} - \Theta^*| &\le \epsilon \iff |(1-\alpha\beta)^T \cdot \Theta^{[0]} - \Theta^*| \le \epsilon\\
  \end{alignat*}

With \(\Theta^*=0\), this simplifies to:
\[
|(1-\alpha\beta)^T \cdot \Theta^{[0]}| \le \epsilon
\]

With logarithms, this can also be expressed as:
\begin{align*}
    |(1-\alpha\beta)^T \cdot \Theta^{[0]}| &\le \epsilon &&\left| : \Theta^{[0]}; \ \log(...)\right.\\
    |(1-\alpha\beta)^T \cdot \Theta^{[0]}| \le \frac{\epsilon}{|\Theta^{[0]}|} &\iff T \ge \frac{\log(\frac{\epsilon}{|\Theta^{[0]}|})}{\log(\frac{1}{|1-\alpha\beta|})} \\
\end{align*}

\vspace{1em}
\textit{The closer $\alpha$ is to 0 or $\frac{2}{\beta}$, the bigger the required number of iterations get. This is due to the logarithm in the denominator getting closer to zero, due to which T increases immensely in size, eventually reaching infinity. This corresponds also with the intuition that, that bigger step sizes overshoot the desired minimum, and take also a lot longer to converge as we hop from "side to side", and lower step sizes just not progressing towards the critical point.}

\end{answer}
